% golden-substrate-whitepaper.tex
% Unified whitepaper: the forced mathematics of the golden substrate,
% the gate field, and the explicit helix.
% Supersedes: HELIX-EXPLICIT-DERIVATION.md v1.0.1,
%             helix_explicit_interactive.html (v1.0.0 companion),
%             ECHO-S-RESEARCH_UNIFIED.md Rev. 2.
\documentclass[11pt]{article}

\usepackage[margin=1in]{geometry}
\usepackage{amsmath,amssymb,amsthm,mathtools}
\usepackage{booktabs,longtable,array}
\usepackage[expansion=false]{microtype}
\usepackage{enumitem}
\usepackage{xcolor}
\usepackage[colorlinks=true,linkcolor=blue!45!black,citecolor=blue!45!black,urlcolor=blue!45!black]{hyperref}

% ---------- epistemic tag macros ----------
\definecolor{cF}{RGB}{20,110,95}   % forced
\definecolor{cD}{RGB}{175,95,10}   % declared
\definecolor{cC}{RGB}{100,100,110} % computed
\definecolor{cI}{RGB}{110,60,150}  % inherited
\definecolor{cE}{RGB}{40,80,160}   % established
\definecolor{cO}{RGB}{170,40,40}   % open
\newcommand{\tagF}{\textcolor{cF}{\textsc{[forced]}}}
\newcommand{\tagD}{\textcolor{cD}{\textsc{[declared]}}}
\newcommand{\tagC}{\textcolor{cC}{\textsc{[computed]}}}
\newcommand{\tagI}{\textcolor{cI}{\textsc{[inherited]}}}
\newcommand{\tagE}{\textcolor{cE}{\textsc{[established]}}}
\newcommand{\tagO}{\textcolor{cO}{\textsc{[open]}}}
\newcommand{\hx}[1]{\textcolor{cC}{\footnotesize\texttt{#1}}}

% ---------- math shorthands ----------
\newcommand{\Q}{\mathbb{Q}}
\newcommand{\Z}{\mathbb{Z}}
\newcommand{\R}{\mathbb{R}}
\newcommand{\C}{\mathbb{C}}
\newcommand{\gp}{\varphi}
\newcommand{\gq}{\psi}
\newcommand{\gapc}{\mathrm{gap}}
\newcommand{\M}{\mathrm{M}}
\newcommand{\Nm}{\mathrm{N}}
\newcommand{\Tr}{\operatorname{Tr}}
\newcommand{\adR}{\operatorname{ad}_R}
\DeclareMathOperator{\lp}{\log^{+}}

\theoremstyle{plain}
\newtheorem{theorem}{Theorem}[section]
\newtheorem{proposition}[theorem]{Proposition}
\newtheorem{lemma}[theorem]{Lemma}
\newtheorem{corollary}[theorem]{Corollary}
\theoremstyle{definition}
\newtheorem{definition}[theorem]{Definition}
\newtheorem{declaration}[theorem]{Declaration}
\theoremstyle{remark}
\newtheorem{remark}[theorem]{Remark}
\newtheorem{guard}[theorem]{Guard}

\title{\bfseries The Golden Substrate, the Gate Field, and the Explicit Helix\\[4pt]
\large A Self-Contained Whitepaper of the Forced Mathematics}
\author{Echo-S-Studios\\[2pt]
\small compiled 2026-07-14; machine-verified (\S\ref{sec:verification})}
\date{Version 1.0 --- July 14, 2026}

\begin{document}
\maketitle

\begin{abstract}
\noindent
This whitepaper consolidates, into one self-contained document, the forced mathematics of
the Echo-S golden substrate: the seed arithmetic of $x^2-x-1$ and of the K-seed
$x^4+5x^2-5$; the keystone matrix $R$ with $R^2=R+I$ and its operator layer
(trace--form duality, the return operator, the $\adR$ spectrum $\{0,\pm\sqrt5\}$);
the two characters of the emission semiring (Mahler measure and $\Z/4\Z$ angle charge)
and the Emission-Gap mechanism excluding Salem numbers from the spectral image;
the trace-down of Lehmer's polynomial and the occupant $\tau_0=\beta+\beta^{-1}$
of the Salem slot; the one-parameter gate family $f_C(x)=C/(1+x)$, its golden multiplier
ladder, and the obstruction theorem (no single real gate carries a helix); and the
explicit helix $\Gamma(\rho)=(r\cos\theta,\ r\sin\theta,\ z)$ built from three declared
atoms, with its gate chart, the rational split gate $C=4/9$ at rapidity exactly $\ln 2$,
the incommensurability theorem $\log_\gp 2\notin\Q$, the field separation
$\Q(\sqrt\gp)\neq\Q(5^{1/4})$, and the registry lattice
($\M(\mathrm{cons})=2\gp^5$ on the $(\ln 2,\ln\gp)$ lattice; $\M(\mathrm{res})$ off it
by a norm obstruction). Every claim carries an epistemic tag; every \tagF{} claim is
proved in this paper or certified by the cited machine checks, and the complete numeric
landmark tables are reproduced. This document supersedes
\texttt{HELIX-EXPLICIT-DERIVATION.md} v1.0.1, its interactive HTML companion (v1.0.0),
and \texttt{ECHO-S-RESEARCH\_UNIFIED.md} Rev.\,2, folding in the corrections identified
by the 2026-07-13 independent soft check.
\end{abstract}

\tableofcontents

\section{Conventions, provenance, and the epistemic ledger}\label{sec:conventions}

\subsection{Epistemic tags}
Every substantive claim in this paper carries one of six tags, used exactly as in the
source corpus:

\begin{center}
\begin{tabular}{@{}lp{12cm}@{}}
\toprule
\tagF & exact consequence of the substrate; proved here and/or machine-verified (check IDs cited)\\
\tagD & a definition or normalization introduced deliberately, stated as such\\
\tagC & numeric display or certified-precision comparison; never load-bearing\\
\tagE & standard external mathematics (textbook or classical literature)\\
\tagI & taken from the source corpus without re-derivation here; provenance cited\\
\tagO & unresolved\\
\bottomrule
\end{tabular}
\end{center}

\noindent
\textbf{Arithmetic discipline.} All decisions are made in exact arithmetic over
$\Q$, $\Q(\sqrt5)$, or (once, trivially) $\Q(\sqrt3)$; statements about
$K=5^{1/4}/\gp$ are decided at the squared level $K^2\in\Q(\sqrt5)$, so that no fourth
root ever crosses a decision boundary. Floating-point values appear only in the
\tagC{} display tables (\S\ref{sec:landmarks}, \S\ref{sec:constants}) and decide nothing.
Machine verification: the shipped harness \texttt{helix\_explicit\_verify.py}
(87 exact checks + 1 computed guard, sympy 1.14.0, exit 0; check IDs \hx{HX01--HX40c})
and the independent soft-check harness \texttt{helix\_softcheck.py}
(89 checks, written cold from the documents' stated claims, exit 0; check IDs
\hx{SC-A1--SC-L2}). See \S\ref{sec:verification}.

\subsection{Provenance and what ``self-contained'' means here}\label{sec:provenance}
This paper replaces three artifacts: the derivation note
\texttt{HELIX-EXPLICIT-DERIVATION.md} v1.0.1 (2026-07-13), its interactive companion
\texttt{helix\_explicit\_interactive.html} (v1.0.0; the $+4$ check-count difference
$83\to 87$ is exactly the v1.0.1 additions \hx{HX39, HX40a--c}), and the corpus
synthesis \texttt{ECHO-S-RESEARCH\_UNIFIED.md} Rev.\,2 covering the fourteen-paper
Echo-S-Research corpus (2026-06-28 to 07-02). Three corrections identified by the
independent soft check are applied in the body of this paper rather than noted as
errata: the minimal-polynomial attribution at rung 1 (\S\ref{sec:fields}), the wording
of the content-polynomial remark (\S\ref{sec:characters}), and the explicit declaration
of the gate iteration (\S\ref{sec:gatefamily}), which the source note used implicitly.

``Self-contained'' means: every \tagF{} statement is either proved in this paper or is
an elementary computation reproduced here and certified by the cited checks. Statements
whose full proofs live in the fourteen corpus papers are stated precisely, tagged
\tagI{}, and scoped exactly as the corpus's own conditionality graph scopes them
(\S\ref{sec:conditionality}); classical results are tagged \tagE{} with canonical
citations. Nothing in this paper depends on the three superseded artifacts.

\subsection{Standing notation}\label{sec:notation}
\[
\gp=\frac{1+\sqrt5}{2},\qquad
\gq=\frac{1-\sqrt5}{2}=-\gp^{-1},\qquad
\tau=\gp^{-1}=\gp-1=\frac{\sqrt5-1}{2},
\]
\[
\gapc=\gp^{-4}=\frac{7-3\sqrt5}{2},\qquad
K^2=1-\gapc=\frac{3\sqrt5-5}{2},\qquad
\beta^2=\frac{5+3\sqrt5}{2}=\gp^2\sqrt5,\qquad
z_c=\frac{\sqrt3}{2}.
\]
$F_n,L_n$ are the Fibonacci and Lucas numbers, extended to all $n\in\Z$ by Binet:
$F_n=(\gp^n-\gq^n)/\sqrt5$, $L_n=\gp^n+\gq^n$. $\M(p)$ denotes the Mahler measure of a
polynomial $p$ (Definition~\ref{def:mahler}); for a matrix or algebra element, the
Mahler-type measure of its spectrum. $\Nm$ is the field norm
$\Q(\sqrt5)\to\Q$, $\Nm(a+b\sqrt5)=a^2-5b^2$.

%=====================================================================
\part{The seed and the keystone}
%=====================================================================

\section{Seed arithmetic}\label{sec:seed}

\subsection{The golden seed}
The seed polynomial is $x^2-x-1$, with roots $\gp,\gq$. The following identities are
elementary and machine-replicated \tagF{} \hx{SC-A7, SC-I16}:
\begin{equation}\label{eq:goldenids}
\gp\gq=-1,\qquad \gp+\gq=1,\qquad \sqrt5=\gp+\gp^{-1}=\gp-\gq,\qquad
\gp^2=\gp+1,\qquad \tau^2+\tau=1 .
\end{equation}
Since $\operatorname{disc}(x^2-x-1)=5$ is squarefree, $\Z[\gp]$ is the maximal order
of $\Q(\sqrt5)$ \tagF{}/\tagE{} \hx{SC-I16}: the golden order is the ring of integers,
not a proper suborder.

\subsection{The gap and the K-seed}\label{sec:kseed}
Define the \emph{gap} $\gapc=\gp^{-4}=(7-3\sqrt5)/2$ and the \emph{K-seed}
\[
p_K(x)=x^4+5x^2-5 .
\]
Substituting $y=x^2$ gives $y^2+5y-5$, whose roots
$y_\pm=(-5\pm3\sqrt5)/2$ straddle zero. The positive root is $K^2=(3\sqrt5-5)/2$,
giving the real pair $\pm K$; the negative root is $-\beta^2$ with
$\beta^2=(5+3\sqrt5)/2=\gp^2\sqrt5$, giving the imaginary pair $\pm i\beta$
\tagF{} \hx{SC-H8}. The seed's basic ledger, all decided at the squared level
\tagF{} \hx{SC-A1--A6, SC-H7}:
\begin{equation}\label{eq:kledger}
\begin{gathered}
K^2=1-\gapc,\qquad
K^4+5K^2-5=0,\qquad
K^4=5\,\gapc,\\
(K\gp)^4=5,\qquad
K^2\beta^2=5,\qquad
\beta^2=K^2\gp^4 .
\end{gathered}
\end{equation}
The last identity reads: the rotation magnitude sits exactly two floors ($\gp^2$ at
the squared level) above the terrain magnitude --- and two is $K$'s own rung index
(\S\ref{sec:rungs}). ``Route 5'' of the K-derivations is the geometric-mean reading
$K^2=\tau(2-\tau)$ \tagF{} \hx{SC-A6}. Irreducibility: $p_K$ is Eisenstein at $5$,
so $[\Q(K):\Q]=4$ and $\Q(K)=\Q(5^{1/4})$ (since $K\gp=5^{1/4}$) \tagF{} \hx{SC-H4}.

The K-seed's charge data --- $\pm K$ on the real (terrain) face at arguments
$\{0,\pi\}$, $\pm i\beta$ on the rotation face at arguments $\{\pm\pi/2\}$, so that
$\chi(\text{K-seed})=\{0,1,2,3\}$ completes $\Z/4\Z$, with Mahler measure
$\M(\text{K-seed})=\beta^2=\gp^2\sqrt5\approx 5.854$ --- is immediate from the root
computation above \tagF{} \hx{SC-H7} (the uniqueness of the K-seed as the catalog's
$\Z/4\Z$-completing straddler, and the even-quartic parity criterion behind it, are
corpus results: paper 13, Props.\ 6.1/6.2 \tagI{}).

The \emph{gap object} is the companion of $x^2-7x+1$, whose large root is $\gp^4$
(the other root is $\gp^{-4}$): $\M(A_{\gapc})=\gp^4$, so
$\ln\M(A_{\gapc})=4\ln\gp$ \tagF{} \hx{SC-E7--E8}. This constant is the full-turn
cost of the helix (\S\ref{sec:winding}).

\section{The keystone matrix and trace--form duality}\label{sec:keystone}

\begin{definition}[Keystone]\label{def:keystone}
$R=\begin{pmatrix}0&1\\ 1&1\end{pmatrix}$, the companion matrix of $x^2-x-1$, so that
\begin{equation}\label{eq:keystone}
R^2=R+I .
\end{equation}
\end{definition}

The corpus derives the keystone rather than positing it: $\gp$ is the smallest Perron
root exceeding $1$ of a $2\times2$ primitive non-negative integer matrix, and each of
the four defining constraints (degree, Perron positivity, growth, integrality) is
load-bearing, with an explicit drop-one witness for each ($\lambda{=}2c$ paper,
Thms.\ 8.3/8.4 and App.\ A) \tagI{}. Everything below needs only \eqref{eq:keystone}.

\begin{proposition}[Fibonacci calculus]\label{prop:fibcalc}\tagF{} \hx{SC-I2--I4}
For all $n\in\Z$ (with $R^{-1}=R-I$):
\[
R^n=F_nR+F_{n-1}I,\qquad
\det R^n=(-1)^n=F_{n+1}F_{n-1}-F_n^2\ \text{(Cassini)},\qquad
\Tr R^n=L_n .
\]
In particular $L_4=\Tr R^4=7$ and the entries of
$R^4=\bigl(\begin{smallmatrix}2&3\\3&5\end{smallmatrix}\bigr)$
are $\{F_3,F_4,F_5\}=\{2,3,5\}$ --- the gate discriminants of the catalog.
\end{proposition}

\begin{proof}
Induction from \eqref{eq:keystone} in both directions; the determinant is
multiplicative with $\det R=-1$. Verified for $n\in[-8,8]$ by \hx{SC-I3, SC-I4}.
\end{proof}

\begin{theorem}[Trace--Form Duality]\label{thm:tfd}\tagF{} \hx{SC-I10--I11, SC-B4}
Let $H=2R-I$. Then $H^2=5I$, the deviation operator collapses,
\[
X_n \coloneqq 2R^n-L_nI = F_nH ,
\]
and for all $n\in\Z$
\begin{equation}\label{eq:tfd}
\tfrac12\Tr(X_n^2)=5F_n^2=L_n^2-4(-1)^n=(\gp^n-\gq^n)^2 .
\end{equation}
Fibonacci supplies the length, Lucas the trace, and Cassini's $\pm4$ is their gap.
\end{theorem}

\begin{proof}
$H^2=4R^2-4R+I=4(R+I)-4R+I=5I$. Then
$X_n=2(F_nR+F_{n-1}I)-L_nI=F_n(2R-I)+(2F_{n-1}+F_n-L_n)I=F_nH$ since
$L_n=F_{n+1}+F_{n-1}=F_n+2F_{n-1}$. Finally
$\tfrac12\Tr(F_n^2H^2)=\tfrac12\,F_n^2\,\Tr(5I)=5F_n^2$, and
$5F_n^2=(\gp^n-\gq^n)^2=L_n^2-4(\gp\gq)^n=L_n^2-4(-1)^n$ by Binet and
\eqref{eq:goldenids}. Verified on $n\in[-8,8]$.
\end{proof}

\section{The operator layer: return, kernel, and the $\sqrt5$ spectrum}\label{sec:operators}

Work in $M_2(\R)\cong\mathrm{Cl}(2,0)$ with basis $\{1,e_1,e_2,J\}$, where
$e_1=\sigma_x=\bigl(\begin{smallmatrix}0&1\\1&0\end{smallmatrix}\bigr)$,
$e_2=\sigma_z=\bigl(\begin{smallmatrix}1&0\\0&-1\end{smallmatrix}\bigr)$,
$e_1^2=e_2^2=1$, $e_1e_2=-e_2e_1$, and $J=e_1e_2$ (so $J^2=-1$; $J$ anticommutes with
every vector). In this basis
\begin{equation}\label{eq:HinClifford}
R=\tfrac12(1+H),\qquad H=2R-I=2e_1-e_2,\qquad H^2=\langle H,H\rangle\,1=5\cdot 1,
\end{equation}
i.e.\ the keystone's traceless part is a Clifford vector of squared length $5$
\tagF{} \hx{SC-I10}.

\begin{proposition}[Sylvester spectra]\label{prop:sylvester}\tagF{} \hx{SC-I5--I7}
On $M_2(\R)$, the Lie self-action $\adR(X)=[R,X]$ has spectrum $\{0,0,+\sqrt5,-\sqrt5\}$;
the Jordan-type map $S(X)=RX+XR$ has spectrum $\{2\gp,\,1,\,1,\,2\gq\}$; and the
\emph{return operator}
\[
L(X)=RX+XR-X=S(X)-X=\tfrac12\{H,X\}
\]
has spectrum $\{+\sqrt5,\,0,\,0,\,-\sqrt5\}$.
\end{proposition}

\begin{proof}
By Sylvester's theorem \tagE{} \cite{HornJohnson1991} the eigenvalues of
$X\mapsto AX+XB$ are the pairwise sums $\alpha_i+\beta_j$ of the spectra of $A$ and
$B^{\!\top}$; with $A=B=R$ (spectrum $\{\gp,\gq\}$) this gives
$\{2\gp,\ \gp+\gq,\ \gp+\gq,\ 2\gq\}=\{2\gp,1,1,2\gq\}$ for $S$, hence
$\{2\gp-1,0,0,2\gq-1\}=\{\sqrt5,0,0,-\sqrt5\}$ for $L=S-\mathrm{id}$, and eigenvalue
\emph{differences} $\{0,0,\pm(\gp-\gq)\}=\{0,0,\pm\sqrt5\}$ for $\adR$. The identity
$L=\tfrac12\{H,\cdot\}$ follows from $R=\tfrac12(1+H)$. All three spectra are
machine-verified as exact $\sqrt5$-expressions.
\end{proof}

$\sqrt5$ therefore appears simultaneously as: the grow eigenvalue of $\adR$; the
$\mathfrak{sl}_2(2,1)$ root length of the Lie layer (primer, Thm.\ 2.5 \tagI{}); the
squared-length of $H$; and (\S\ref{sec:exchange}) the forced exchange rate $\lambda$.

\begin{theorem}[Kernel of the return operator]\label{thm:kernel}\tagF{} \hx{SC-I7--I9}
$\ker L$ is exactly two-dimensional:
\[
\ker L=\operatorname{span}\{\,e_1+2e_2,\ J\,\}.
\]
Moreover, on Clifford vectors $v$, $L$ acts as the scalar pairing
$L(v)=\langle H,v\rangle\,1$; on the scalar--$H$ plane it acts as the off-diagonal
block $\bigl(\begin{smallmatrix}0&5\\1&0\end{smallmatrix}\bigr)$ with eigenvalues
$\pm\sqrt5$.
\end{theorem}

\begin{proof}
For vectors $u,v$, $\{u,v\}=2\langle u,v\rangle\,1$ in $\mathrm{Cl}(2,0)$; with
$H=2e_1-e_2$, $L(ae_1+be_2)=\langle H,\,ae_1+be_2\rangle\,1=(2a-b)\,1$, which vanishes
iff $b=2a$, i.e.\ on $\operatorname{span}\{e_1+2e_2\}$. The pseudoscalar $J$
anticommutes with every vector, so $\{H,J\}=0$ and $L(J)=0$. Finally $L(1)=H$ and
$L(H)=H^2=5\cdot1$, giving the $\pm\sqrt5$ plane. Dimensions: $2+2=4$. The kernel
dimension and basis are machine-verified by exact nullspace computation
(\hx{SC-I7, SC-I8}), matching the corpus's rational-elimination result
(Residual Return, \S6 \tagI{}).
\end{proof}

\begin{remark}[The disproved design hint]\label{rem:64}\tagF{} \hx{SC-I9}
The corrected corpus Remark 6.4 pins a refuted design note: the vector $-e_1+e_2$ is
\emph{not} in $\ker L$; indeed by Theorem~\ref{thm:kernel},
$L(-e_1+e_2)=\langle H,\,-e_1+e_2\rangle\,1=(-2-1)\,1=-3\cdot 1$, replicated here
exactly. The scalar formula makes the disproof transparent: among vectors, the kernel
is precisely $H^{\perp}$, and $-e_1+e_2$ is not orthogonal to $H$. (The convention
matters: only the companion orientation $R=\bigl(\begin{smallmatrix}0&1\\1&1
\end{smallmatrix}\bigr)$, $e_1=\sigma_x$, $e_2=\sigma_z$ makes both the kernel basis
and the $-3$ land; this paper fixes that convention once, in
\eqref{eq:HinClifford}.)
\end{remark}

For completeness of the representation layer \tagI{}/\tagE{} \hx{SC-I15}: the coupled
$\mathfrak{sl}_2$ targets have $\dim V_m=m+1$ and Casimir $\tfrac12 m(m+2)$, giving
$15/2$ on $V_3$ and $12$ on $V_4$; squaring on the semiring is the Adams operation
$\psi^2$ (Operator Algebra paper, Thm.\ 3.2 \tagI{}).

\section{The two characters and the Emission-Gap mechanism}\label{sec:characters}

\begin{definition}[Mahler measure]\label{def:mahler}\tagE{} \cite{Lehmer1933,EverestWard1999}
For a monic $p(x)=\prod_i(x-\alpha_i)\in\Z[x]$,
$\M(p)=\prod_i\max(1,|\alpha_i|)$, i.e.\ $\log\M(p)=\sum_i\lp|\alpha_i|$.
An algebraic integer $\beta>1$ is a \emph{Salem number} if it is reciprocal
(conjugate-set closed under $\alpha\mapsto1/\alpha$) with all conjugates in
$\overline{\mathbb D}$ and at least one on the unit circle \tagE{} \cite{Salem1963};
a \emph{Pisot number} if all conjugates other than $\beta$ lie strictly inside the
circle \tagE{} \cite{Siegel1944}. Lehmer's polynomial
$\ell(x)=x^{10}+x^9-x^7-x^6-x^5-x^4-x^3+x+1$ has
$\M(\ell)=1.1762808\ldots$, the smallest known measure $>1$; whether measures
accumulate at $1$ is Lehmer's problem \tagO{}, external \cite{Lehmer1933}.
Smyth's theorem \tagE{} \cite{Smyth1971}: a nonreciprocal algebraic integer $\alpha\neq0$
has $\M(\alpha)\ge\mu_S=1.3247179\ldots$, the plastic number (root of $x^3-x-1$),
which is also the smallest Pisot number \tagE{} \cite{Siegel1944}.
\end{definition}

\subsection{Catalog, operators, and the two characters}
The corpus catalog consists of the companion matrices of the seeds
\[
\{\ \gp,\ \tau,\ \sqrt2,\ \sqrt3,\ \sqrt5,\ \ \gapc\ (x^2{-}7x{+}1),\ \ K\ (x^4{+}5x^2{-}5)\ \};
\]
the \emph{spectral operators} are those whose output spectrum is forced by input
spectra --- $\otimes$ (Kronecker product: spectrum $\{\alpha_i\beta_j\}$
\tagE{} \cite{HornJohnson1991}), $\oplus$ (direct sum: union), squaring
($=$ Adams $\psi^2$), minimal polynomial, and cyclotomic factors $\Phi$. The
\emph{spectral image} $S$ is the closure of the catalog under these. Two invariants
turn out to be the two characters of the resulting $\lambda$-ring
(Operator Algebra, Thm.\ 8.1 \tagI{}):

\begin{theorem}[Tropical tensor law for $\M$]\label{thm:tropical}\tagF{} \hx{SC-I14}
On spectra, $\log\M$ is additive on $\oplus$ and \emph{tropical} on $\otimes$:
\[
\log\M(A\otimes B)=\sum_{i,j}\bigl(\log|\alpha_i|+\log|\beta_j|\bigr)^{+}.
\]
Witness: $\gp\otimes\gp$ has spectrum $\{\gp^2,\ \gp\gq,\ \gq\gp,\ \gq^2\}
=\{\gp^2,-1,-1,\gq^2\}$, so $\M=\gp^2$ --- \emph{not} $\gp^4$; the naive product law
$\M(A)^{\deg B}\M(B)^{\deg A}$ agrees with the tropical law only off the unit circle
(this is the corpus's v4 correction, pinned by exactly this witness).
\end{theorem}

\begin{proof}
Spectrum of $A\otimes B$ is $\{\alpha_i\beta_j\}$, so
$\log\M(A\otimes B)=\sum_{i,j}\lp|\alpha_i\beta_j|
=\sum_{i,j}(\log|\alpha_i|+\log|\beta_j|)^{+}$. For the witness, $|\gp\gq|=1$ exactly
by \eqref{eq:goldenids}, so the two cross terms contribute $(0)^{+}=0$ and
$|\gq^2|<1$ contributes $0$; only $(2\log\gp)^{+}=2\log\gp$ survives.
\end{proof}

The second character is the \emph{angle charge} $\chi$: when every conjugate argument
of an object lies on the rational-angle lattice $\tfrac{2\pi}{n}\Z$, the object carries
charge group $\Z/n\Z$; $\chi$ adds (as a sumset) under $\otimes$, doubles under
$\psi^2$, and unions under $\oplus$ \tagI{} (Charge--Measure Coupling, Thms.\ 3.1--4.4).
For the emission image the lattice is $(\pi/2)\Z$ and the charge is $\Z/4\Z$.

\begin{remark}[Content polynomial, wording fixed]\label{rem:content}
$x^4-1=\Phi_1\Phi_2\Phi_4$ \tagF{} \hx{SC-I13}. The generative-emptiness slogan
attached to this factorization is stated here in its unambiguous form: the
\emph{charge} carries the full $\Z/4\Z$ (all four arguments $0,\pi/2,\pi,3\pi/2$
occur), while the \emph{real terrain face} of the root set is only
$\{\pm1\}\cong\Z/2\Z$; the pair $\pm i$ lives in the rotation fiber. (The superseded
summary's compressed phrasing ``only $\Z/2\Z$ on the circle'' read falsely at face
value, since all four fourth roots of unity lie on the unit circle; the intended
terrain/fiber split is the statement above.)
\end{remark}

\subsection{The Emission-Gap mechanism}\label{sec:emissiongap}
The corpus's central confinement result is proved by \emph{excision}, not by bounding
all integer polynomials. The mechanism has four steps; the first three are elementary
and reproduced here, the packaging over the full operator closure is the corpus
theorem.

\begin{lemma}[Angle confinement of the catalog]\label{lem:angles}\tagF{}
Every eigenvalue of every catalog seed has argument in $(\pi/2)\Z$.
\end{lemma}
\begin{proof}
$\gp,\tau,\sqrt2,\sqrt3,\sqrt5$ and the gap object are totally real: arguments in
$\{0,\pi\}$. The K-seed's spectrum is $\{\pm K,\pm i\beta\}$ (\S\ref{sec:kseed}):
arguments $\{0,\pi,\pm\pi/2\}$.
\end{proof}

\begin{lemma}[Operator preservation]\label{lem:preserve}\tagF{}
The spectral operators preserve the lattice $\arg\in(\pi/2)\Z$: $\otimes$ adds
arguments, $\psi^2$ doubles them, $\oplus$ unions them, and minimal-polynomial/$\Phi$
extraction selects subsets. (The lattice is closed under addition and doubling ---
\emph{not} under halving; no operator halves an angle.)
\end{lemma}

\begin{lemma}[On-circle emissions are fourth roots of unity]\label{lem:mu4}\tagF{}
If $\lambda$ is an eigenvalue of an element of the spectral image $S$ with
$|\lambda|=1$, then $\lambda\in\{\pm1,\pm i\}=\mu_4$.
\end{lemma}
\begin{proof}
Immediate from Lemmas \ref{lem:angles}--\ref{lem:preserve}: $|\lambda|=1$ and
$\arg\lambda\in(\pi/2)\Z$.
\end{proof}

\begin{lemma}[Salem circle-conjugates are never torsion]\label{lem:salemtorsion}\tagF{}/\tagE{}
No conjugate of a Salem number that lies on the unit circle is a root of unity.
\end{lemma}
\begin{proof}
If a root of unity $\zeta$ were a conjugate of the Salem number $\beta$, the minimal
polynomial of $\beta$ would equal the minimal polynomial of $\zeta$, a cyclotomic
polynomial, whose roots all have modulus $1$ --- contradicting $\beta>1$.
\end{proof}

\begin{theorem}[Emission-Gap, scoped]\label{thm:emissiongap}
Over the spectral image $S$: no element of $S$ has a Salem spectrum, and
$\M(S)\subseteq\{1\}\cup[\mu_S,\infty)$, with the floor realized inside the image at
$\gp$. \tagF{} over $S$ (Lemmas \ref{lem:angles}--\ref{lem:salemtorsion} $+$ Smyth
\tagE{} $+$ the empty quadratic strip $(1,\gp)$); the entry-level closures
(circulant \tagF{}; commutator in-context \tagF{}, uniform via
$\adR$-derivation and the signature lattice of $\Q(\sqrt2,\sqrt3,5^{1/4})$ \tagF{};
unbounded free commutator \tagO{}) and the universal statement over all
charge-admissible objects (\tagF{} for degree $\le2$; \tagC{}/plausible in general)
are corpus results (Emission-Gap paper Thm.\ 3.2, Lehmer's Box paper Thms.\ 5.3--5.4)
\tagI{}. Equivalent forms: Mahler, entropy ($h=\log\M$
\tagE{} \cite{LindSchmidtWard1990}), spectral (no on-circle-expanding channel), and
metric (no emitted field places a complex embedding on the circle) \tagI{}.
\end{theorem}

The one door of Lehmer's Box --- the free commutator at unbounded size, through which
the full Salem spectrum (Lehmer's number included) re-enters --- remains \tagO{};
the corpus closes it structurally where the framework's commutator is $\adR$ (a
derivation with difference spectrum $\{0,\pm\sqrt5\}$, Prop.~\ref{prop:sylvester}) and
locks it operationally elsewhere ($\beta<\gp$ decided by the sign of an element of
$\Q(\sqrt5)$) \tagI{}. The corpus's own closing scope statement is retained verbatim
in intent: \emph{``We have not shown they cannot come close to $1$; we have shown the
construction never goes to fetch them.''}

\section{Trace-down, the Salem slot, and the parity floor}\label{sec:tracedown}

\subsection{The trace-down map and Lehmer's instance}
For a reciprocal polynomial of degree $2d$, the substitution $t=x+x^{-1}$ produces a
degree-$d$ \emph{trace-down} $T(t)$ with $x^d\,T(x+x^{-1})=p(x)$. Downstairs, the three
channels are the three regions of $\R$ cut by $t=\pm2$: on-circle conjugates land in
$(-2,2)$, expanding ones in $(2,\infty)$. A Salem number is exactly a
\emph{flip-straddler}: one trace-conjugate in $(2,\infty)$, the rest in $(-2,2)$ ---
so ``Salem'' is an upstairs artifact of the $2{:}1$ lift, and the slot's occupant is
the totally real Perron number $\tau_0=\beta+\beta^{-1}$ \tagI{} (Salem Slot paper).

\begin{proposition}[Lehmer's trace-down, exact]\label{prop:lehmertrace}\tagF{} \hx{SC-J3--J5}
With $\ell(x)=x^{10}+x^9-x^7-x^6-x^5-x^4-x^3+x+1$ and
$p_k \coloneqq x^k+x^{-k}$ expressed by the Chebyshev-type reductions
$p_1=t$, $p_2=t^2-2$, $p_4=t^4-4t^2+2$, $p_5=t^5-5t^3+5t$:
\[
\frac{\ell(x)}{x^5}=p_5+p_4-p_2-p_1-1
=t^5+t^4-5t^3-5t^2+4t+3 \eqqcolon T(t),
\]
so $x^5\,T(x+x^{-1})=\ell(x)$ exactly. Its largest root is
$\tau_0=\beta+\beta^{-1}=2.0264179\ldots$ where $\beta=\M(\ell)=1.1762808\ldots$;
$\M(T)=5.615601\ldots$ \tagC{}.
\end{proposition}

\begin{proof}
Expand $\ell/x^5=(x^5+x^{-5})+(x^4+x^{-4})-(x^2+x^{-2})-(x+x^{-1})-1$ and substitute
the reductions; the polynomial identity is machine-verified exactly (\hx{SC-J3}), the
root and measure values at 60 decimal places (\hx{SC-J1, SC-J4, SC-J5}).
\end{proof}

Three structural facts about the slot, elementary at this level \tagF{}/\tagI{}:
the branch identity $\tau_0-2=(\beta-1)^2/\beta$ (a Salem of height $1+\delta$
redirects only $O(\delta^2)$ past the flip); the interval $(2,\sqrt5\,]$ is the trace
image of real pairs $\beta\in(1,\gp]$, with $t=\sqrt5$ lifting to $\{\gp,\gp^{-1}\}$
--- so the redirections converge to $\sqrt5=\gp+\gp^{-1}$, the $\adR$ grow eigenvalue
(Prop.~\ref{prop:sylvester}), with linear slope $\gp^{-1}$ and geometric ratio $\gp$
along the canonical family \tagI{}; and the accumulation points of redirections are
the Pisot traces, least $\sqrt5$, next the plastic
$\mu_S+\mu_S^{-1}=2.0795956\ldots$ \tagC{} \hx{SC-J9}. The identity making the trace
map canonical rather than imported is $\gp\gq=-1$ (Salem Slot, Prop.\ 7.1 \tagI{}).

\subsection{The parity-graded floor and the $q_k$ construction}\label{sec:parity}

\begin{proposition}[$\mu(\text{even})\le\gp$, constructive]\label{prop:qk}\tagF{} \hx{SC-J11}
For every $k\ge1$, $q_k(x)=x^{2k}+x^k-1$ has $\M(q_k)=\gp$.
\end{proposition}

\begin{proof}
Substituting $y=x^k$ gives $y^2+y-1$ with roots $\tau$ and $-\gp$. The $k$ roots with
$x^k=-\gp$ each have modulus $\gp^{1/k}$, contributing $(\gp^{1/k})^k=\gp$; the $k$
roots with $x^k=\tau$ lie strictly inside the circle. (That $q_k$ carries charge
$\Z/2k\Z$ with all charges attained, and the pentagon-free character of the
substitution, are the corpus's parity-floor packaging \tagI{}.) Verified numerically
for $k\le4$ at 60 dps.
\end{proof}

The floor structure, with the corpus's honest strength split \tagI{} (Charge--Measure
Coupling v4, Thm.\ 5.1; Z/5Z paper): $\mu(\text{even})\le\gp$ \tagF{} (above);
$\mu(\text{even})\ge\gp$ plausible; $\mu(\text{odd})\ge\mu_S$ \tagF{} (Smyth, sharpened
at $\Z/5\Z$: non-reciprocal charge-$\Z/5\Z$ $\Rightarrow \M\ge\mu_S$; reciprocal
$\Rightarrow \M\in\{1\}\cup[\gp^2,\infty)$; pure pentagon quartics
$\M\in\{1\}\cup[\gp^4,\infty)$); $\mu(\text{odd})=2$ \tagC{}, attained at $x^5-2$
(all five roots of modulus $2^{1/5}=1.1487\ldots$, so $\M=2$ exactly \tagF{}
\hx{SC-J10}), with the residual named exactly: non-reciprocal charge-$\Z/5\Z$ objects
in $[\mu_S,2)$ \tagO{}. The mechanism of the parity coupling is the $\pi$-ray
obstruction: $\gq$ sits at argument $\pi$, and $\pi\in\tfrac{2\pi}{n}\Z$ iff $n$ is
even \tagI{}.

\section{The exchange rate $\lambda=2c$ and the forced scalar $\sqrt5$}\label{sec:exchange}

The Residual-Return machine treats a number field $\Q(\theta)$ as an exact learning
geometry: the trace form $G=M^{\top}M$ (Minkowski embedding $M$; $\det G=d_K$ up to
index squared) prices a residual $r$, and growth adjoins a seed $\theta$ iff
$\|r\|_G^2\ge\lambda\log\M(\theta)$. Reading $\|r\|_G^2$ as a second-order
Kullback--Leibler divergence with Fisher metric $cG$ gives
$D_{\mathrm{KL}}\approx\tfrac1{2c}\|r\|_G^2$, and the MDL balance
$D_{\mathrm{KL}}\ge\log\M$ yields the threshold coefficient
\begin{equation}\label{eq:lambda2c}
\boxed{\ \lambda=2c\ }\qquad\text{for every } c>0 \quad\tagI{}\ (\lambda{=}2c\text{ paper, Thm.\ 2.2}).
\end{equation}
By \v{C}encov's theorem \tagE{} \cite{Cencov1982} no invariance argument can remove
the conformal scale $c$: the honest statement is ``$\lambda=2c$ is forced; $c$ is a
declared positive scale,'' never ``zero free parameters'' \tagD{}. The framework's own
structure then pins the canonical instance: $C=1$ is simultaneously the minimal
structure realizing the ternary growth decision and the terminal firewall image fixed
by the keystone --- two premise-disjoint derivations converging on
\eqref{eq:keystone} --- giving
\[
C=1\ \Rightarrow\ c=\tfrac{\sqrt5}{2}\ \Rightarrow\ \lambda=\sqrt5
\qquad\tagI{}
\]
($\lambda{=}2c$ paper, Thm.\ 7.7 --- forced \emph{given} the two constitutive
axioms: ternary arity and the keystone).
The forced-canonical floor is therefore $\lambda\log\mu_S=\sqrt5\,\log\mu_S=
0.628781\ldots$ \tagC{} \hx{SC-J13}; the shipped probe values $2c\log\mu_S$ at
$c\in\{1,4,8\}$ are $0.5624$, $2.2496$, $4.4992$ \tagC{} \hx{SC-J12}. On the helix
chart (\S\ref{sec:chart}), $\sqrt5$ reappears as the rung-1 exchange rate
$\lambda(z_1)=\sqrt5$ --- the same scalar, met from the geometric side.

%=====================================================================
\part{The gate field}
%=====================================================================

\section{The gate family, declared}\label{sec:gatefamily}

The source note used the gate family through its fixed-point polynomial
$g_C(x)=x^2+x-C$ without printing the iteration; this paper declares it once. (The
declaration is pinned, not free: it is the unique reading under which \emph{all} of
the corpus identities used below --- $m_+m_-=1$, $m+\tfrac1m=-\tfrac1C-2$,
$m_1=-\gp^{-2}$, the lens multipliers $\{-\tfrac14,-4\}$, $\pm i$ at $C=-\tfrac12$,
and the pentagon quanta --- hold simultaneously.)

\begin{declaration}[Gate iteration]\label{dec:gate}\tagD{}
For $C\in\R$, the gate $G_C$ is the iteration
\[
f_C(x)=\frac{C}{1+x},
\]
whose fixed points solve $g_C(x)=x^2+x-C=0$, i.e.\
$u_\pm=\tfrac12(-1\pm\sqrt D)$ with discriminant $D=1+4C$, and whose multiplier at a
fixed point $u$ is
\[
m=f_C'(u)=-\frac{C}{(1+u)^2}.
\]
The \emph{flip} is the sign change of $D$ at $C=-\tfrac14$ (real
$\leftrightarrow$ complex fixed points; terrain $\leftrightarrow$ rotation face).
\end{declaration}

\begin{proposition}[Flip reciprocity and the trace identity]\label{prop:reciprocity}\tagF{} \hx{SC-F2--F3}
For every $C\neq0$ the two multipliers satisfy
\[
m_+m_-=1,\qquad\qquad m+\frac1m=-\frac1C-2\in\R .
\]
\end{proposition}

\begin{proof}
By Vieta on $g_C$: $u_++u_-=-1$ and $u_+u_-=-C$, so
$(1+u_+)(1+u_-)=1-1-C=-C$, hence
$m_+m_-=C^2/\bigl((1+u_+)(1+u_-)\bigr)^2=C^2/C^2=1$. For the trace:
$m_++m_-=-C\bigl[(1+u_+)^{-2}+(1+u_-)^{-2}\bigr]
=-C\,\frac{(1+u_+)^2+(1+u_-)^2}{C^2}$; and
$(1+u_+)^2+(1+u_-)^2=\bigl((1+u_+)+(1+u_-)\bigr)^2-2(1+u_+)(1+u_-)=1+2C$, giving
$m_++m_-=-(1+2C)/C=-\tfrac1C-2$. Since $m_-=1/m_+$ this is $m+\tfrac1m$. Both
identities are machine-verified for generic symbolic $C$.
\end{proof}

\section{The golden multiplier ladder}\label{sec:ladder}

\begin{theorem}[Ladder in Lucas--Fibonacci dress]\label{thm:ladder}\tagF{} \hx{SC-B1--B3}
For $n\ge1$ set $C_n=\dfrac{1}{(\gp^{\,n}-\gp^{-n})^2}$. Then the attracting fixed
point, multiplier, rapidity, discriminant, and exchange rate of $G_{C_n}$ are
\[
u_n=\frac{1}{\gp^{2n}-1},\qquad
m_n=-\gp^{-2n},\qquad
\rho_n=n\ln\gp,
\]
\[
\sqrt{D_n}=\frac{\gp^{\,n}+\gp^{-n}}{\gp^{\,n}-\gp^{-n}},\qquad
\lambda_n=\gp^{2n}-\gp^{-2n}=\sqrt5\,F_{2n},
\]
with the Trace--Form Duality split (Theorem~\ref{thm:tfd})
\[
C_n=\begin{cases}1/L_n^2,& n\ \text{odd},\\[2pt] 1/(5F_n^2),& n\ \text{even},\end{cases}
\qquad
\sqrt{D_n}=\begin{cases}\sqrt5\,F_n/L_n,& n\ \text{odd},\\[2pt] L_n/(\sqrt5\,F_n),& n\ \text{even}.\end{cases}
\]
In particular $C_1=1$ (the \emph{golden gate}) and $C_2=\tfrac15$ (the \emph{K-gate});
explicitly $C_3=\tfrac1{16}$, $C_4=\tfrac1{45}$, $C_5=\tfrac1{121}$.
\end{theorem}

\begin{proof}
$u_n$ solves $g_{C_n}$: with $w=\gp^{2n}$,
$u_n^2+u_n=\frac{1+(w-1)}{(w-1)^2}=\frac{w}{(w-1)^2}=\frac{1}{(\gp^n-\gp^{-n})^2}=C_n$.
Multiplier: $1+u_n=\frac{w}{w-1}$, so
$m_n=-C_n\frac{(w-1)^2}{w^2}=-\frac{w}{(w-1)^2}\cdot\frac{(w-1)^2}{w^2}=-w^{-1}
=-\gp^{-2n}$, whence $|m_n|=e^{-2\rho_n}$ with $\rho_n=n\ln\gp$. Discriminant:
$D_n=1+4C_n=\frac{(\gp^n-\gp^{-n})^2+4}{(\gp^n-\gp^{-n})^2}
=\Bigl(\frac{\gp^n+\gp^{-n}}{\gp^n-\gp^{-n}}\Bigr)^2$. The parity dress follows from
$\gp^{-n}=(-1)^n\gq^{\,n}$: for $n$ odd, $\gp^n-\gp^{-n}=\gp^n+\gq^n=L_n$ and
$\gp^n+\gp^{-n}=\sqrt5F_n$; parities swap for $n$ even. Finally
$\lambda_n$ is the chart exchange rate at the rung (Theorem~\ref{thm:chart} below):
$\lambda_n=\frac1{1-z_n^2}-(1-z_n^2)=\gp^{2n}-\gp^{-2n}=\sqrt5F_{2n}$ by Binet.
All of it is machine-verified for $n=1..6$ and, for the TFD identity, on
$n\in[-8,8]$.
\end{proof}

\section{Special gates}\label{sec:specialgates}

\begin{proposition}[The K-gate certificate]\label{prop:kgate}\tagF{} \hx{SC-F1, SC-C4}
At the K-gate $C_2=\tfrac15$: the occupation is the seed's inverse Mahler measure,
\[
u_2=\frac{1}{\gp^4-1}=\frac{1}{\beta^2}=\frac{1}{\M(\text{K-seed})},
\]
and the multiplier is minus the gap, $m_2=-\gp^{-4}=-\gapc$. The gap paper's
self-consistency derivative $f'(K)=-\gapc$ is certified by the exact congruence
\[
5y\ \equiv\ (1-y)^2(5+y)^3 \pmod{\,y^2+5y-5\,},
\]
an identity in $\Q[y]$ evaluated at the K-seed spectrum $y\in\{K^2,-\beta^2\}$.
\end{proposition}

\begin{proof}
$\gp^4-1=\gp^2(\gp^2-\gp^{-2})=\gp^2\sqrt5=\beta^2$ by \eqref{eq:goldenids} and
\eqref{eq:kledger}. The congruence, reduced by $y^2\equiv 5-5y$:
$(1-y)^2\equiv 6-7y$; $(5+y)^2\equiv 30+5y$; $(5+y)^3\equiv175+30y$; and
$(6-7y)(175+30y)=1050-1045y-210y^2\equiv1050-1045y-210(5-5y)=5y$. Machine-verified as
an exact polynomial remainder (\hx{SC-F1}); the interpretation of the left side as
$5\,|f'(K)|$-data along the K-route is the gap paper's \tagI{}, but on the chart of
\S\ref{sec:chart} the same value is forced independently: $m(z_2)=-(1-K^2)=-\gapc$.
\end{proof}

\begin{proposition}[Rotation gates and the quanta inventory]\label{prop:quanta}\tagF{} \hx{SC-F5--F7}
Past the flip ($C<-\tfrac14$) the multipliers are unimodular. Three rotation quanta
exist in the substrate:
\begin{itemize}[nosep]
\item $\pi$: every terrain multiplier is negative (each step crosses the fixed point
--- a half-turn on the line);
\item $\pi/2$: at $C=-\tfrac12$ the multipliers are exactly $\pm i$ (the odd-charge
generators of $\Z/4\Z$);
\item $2\pi/5,\,4\pi/5$: at $C=-\gp^{-2}$ and $C=-\gp^{2}$ the multiplier traces are
$m+\tfrac1m=\gp^2-2=\tau=2\cos72^\circ$ and $\gp^{-2}-2=-\gp=2\cos144^\circ$
respectively --- primitive fifth roots of unity.
\end{itemize}
\end{proposition}

\begin{proof}
At $C=-\tfrac12$ the fixed points are $u=\tfrac12(-1\pm i)$, so
$1+u=\tfrac12(1\pm i)$ and $(1+u)^2=\pm\tfrac i2$, giving
$m=-C/(1+u)^2=\tfrac12\cdot\tfrac{2}{\pm i}=\mp i$; $|m|=1$, arguments $\pm\pi/2$. For the pentagon gates apply
Proposition~\ref{prop:reciprocity}: $m+\tfrac1m=-\tfrac1C-2$ gives
$\gp^2-2=\tau$ at $C=-\gp^{-2}$ and $\gp^{-2}-2=-\gp$ at $C=-\gp^{2}$, i.e.\
$2\cos(2\pi/5)$ and $2\cos(4\pi/5)$, so $m$ is a primitive fifth root of unity. All
three are machine-verified exactly.
\end{proof}

\section{The obstruction theorem}\label{sec:obstruction}

\begin{theorem}[No single-gate helix]\label{thm:obstruction}\tagF{} \hx{SC-F2--F4}
Let $C\in\R$ and let $m$ be a fixed-point multiplier of $G_C$. Then $m\in\R$ or
$|m|=1$. In particular no real gate level linearizes to a genuine spiral
($0<|m|<1$ with $m\notin\R$).
\end{theorem}

\begin{proof}
By Proposition~\ref{prop:reciprocity} the partner of $m$ is $1/m$ and
$m+\tfrac1m=-\tfrac1C-2\in\R$. Writing $m=s\,e^{i\vartheta}$ with $s>0$:
$\operatorname{Im}\bigl(m+\tfrac1m\bigr)=(s-\tfrac1s)\sin\vartheta$; its vanishing
forces $s=1$ or $\sin\vartheta=0$.
\end{proof}

\begin{corollary}[Helicity needs the straddler]\label{cor:straddler}\tagF{}$+$\tagI{}
Simultaneous rotation and contraction cannot be carried by any single real gate; it
requires composing the two faces of the flip. The K-seed $x^4+5x^2-5$ is the unique
catalog seed straddling the flip (real pair $\pm K$ on the terrain face; imaginary
pair $\pm i\beta$ supplying the quarter-turn; uniqueness and the even-quartic parity
criterion: paper 13, Props.\ 5.1/6.1--6.2, Thm.\ 7.1 \tagI{}). Its two faces multiply
to the field prime, $K^2\beta^2=5$, with $\M(\text{K-seed})=\beta^2=\gp^2\sqrt5$ and
$\beta^2=K^2\gp^4$ \tagF{} \eqref{eq:kledger} --- the rotation magnitude exactly two
floors above the terrain magnitude, $K$'s own rung index.
\end{corollary}

This sharpens the corpus's coupling theorem (``a helix cannot couple to a copy of
itself, only to the rotation face of its own flip'') to the gate family: the flip
identity of Proposition~\ref{prop:reciprocity} \emph{is} the obstruction, and the
straddler is the minimal way around it.

%=====================================================================
\part{The explicit helix}
%=====================================================================

\section{Inherited data and the three declared atoms}\label{sec:atoms}

The orthogonal-partner paper dissolves the legacy L4 helix as a load-bearing object;
what survives its harness, both \tagF{}, is the pair \emph{winding $=$ charge
$\chi\in\Z/4\Z$} (the quarter-turn supplied by the K-seed's imaginary pair $\pm i\beta$)
and \emph{pitch $=$ Mahler magnitude $\M\ge\gp$} (terrain growth in floors of
$\ln\gp$). This Part draws the survivors as an explicit curve
\[
\Gamma(\rho)=\bigl(r(\rho)\cos\theta(\rho),\ r(\rho)\sin\theta(\rho),\ z(\rho)\bigr),
\qquad \rho\in[0,\infty),
\]
from exactly three declared atoms plus one inherited profile. No $0.1$-type threshold
appears anywhere; the gap enters only through forced identities.

\smallskip\noindent\textbf{Inherited} \tagI{}: the seed arithmetic and ladder of
\S\S\ref{sec:seed}--\ref{sec:specialgates} (forced in source, re-proved here), and the
radius profile
\begin{equation}\label{eq:radius}
r(z)=K\sqrt{z/z_c}\ \ \text{for } z\in[0,z_c],\qquad r(z)=K\ \ \text{for } z\in[z_c,1],
\qquad z_c=\tfrac{\sqrt3}{2},
\end{equation}
which is read off the L4 source (K-VARIABLE-DERIVATIONS \S8) and \emph{not} derived
here (\tagO{}~1, \S\ref{sec:open}). Its internal consequences below (unique $z=r$
point, kink) are \tagF{} given the profile.

\begin{declaration}[D1: the chart]\label{dec:d1}\tagD{}
The vertical coordinate is calibrated by \emph{co-intensity $=$ local multiplier
magnitude}:
\begin{equation}\label{eq:d1}
1-z^2=e^{-2\rho}=|m|(\rho),\qquad\text{equivalently}\qquad
z(\rho)=\sqrt{1-e^{-2\rho}},\quad \rho(z)=-\tfrac12\ln(1-z^2).
\end{equation}
Justification, two-fold: (i) it extends the gap paper's point identity
$|f'(K)|=\gapc=1-K^2$ from a single height into the defining property of the
coordinate --- the chart does to every height exactly what the main theorem does to
$K$; (ii) it is consistency-checked against the ladder: the chart occupation
$u(z)=(1-z^2)/z^2$ pulls back to the ladder's own fixed points, $u(z_n)=u_n$, at every
rung (Theorem~\ref{thm:chart}).
\end{declaration}

\begin{declaration}[D2: selection principle]\label{dec:d2}\tagD{}
\emph{The winding realizes the seed's charge completion.} Given D2, the angular
quantum per rung is \tagF{} a generator of $\Z/4\Z$, i.e.\ $\chi=\pm\pi/2$ --- the
multipliers $\pm i$ of the $C=-\tfrac12$ gate (Prop.~\ref{prop:quanta}), exactly the
charge the K-seed supplies. Orientation $+i$ (counterclockwise) is a convention. The
competing corpus quanta --- $\pi$ (terrain sign flip) and $2\pi/5$ (pentagon gates)
--- generate only $\Z/2\Z$ and $\Z/5\Z$ and are excluded \emph{by the principle}; the
quanta themselves are \tagF{} corpus data (Prop.~\ref{prop:quanta}).
\end{declaration}

\begin{declaration}[D3: rate normalization]\label{dec:d3}\tagD{}
One charge quantum per entropy floor: $\dfrac{d\theta}{d\rho}=\dfrac{\pi/2}{\ln\gp}$.
This is the minimal coupling of the two surviving \tagF{} anatomical facts (winding
$\in\Z/4\Z$; pitch in floors of $\ln\gp$). Any other constant rate rescales the pitch
without changing the qualitative structure (Prop.~\ref{prop:invariance}).
\end{declaration}

Everything in this Part is \tagF{} given the inherited data $+$ D1--D3, except where
tagged otherwise.

\section{The vertical chart: every height is a gate}\label{sec:chart}

\begin{theorem}[The gate chart]\label{thm:chart}\tagF{} \hx{SC-C1--C8}
Under D1, with $z\in(0,1]$, the gate data become functions of height:
\[
u(z)=\frac{1-z^2}{z^2},\qquad
m(z)=-(1-z^2),\qquad
C(z)=\frac{1-z^2}{z^4},
\]
\[
\sqrt{D(z)}=\frac{2-z^2}{z^2},\qquad
\lambda(z)=\frac{z^2(2-z^2)}{1-z^2},
\]
and:
\begin{enumerate}[nosep,label=(\alph*)]
\item $u(z)$ is a fixed point of $G_{C(z)}$ with multiplier exactly $m(z)$, so
$|m|=1-z^2$ holds \emph{chart-wide}, not only at rungs;
\item $1+4C(z)=\bigl((2-z^2)/z^2\bigr)^2$ exactly;
\item $C(z)$ is a strictly decreasing bijection $(0,1]\to[0,\infty)$
($C'(z)=2(z^2-2)/z^5<0$ on the interval; $C(1)=0$, $C\to\infty$ as $z\to0^+$): every
attracting gate level occurs at exactly one height. The rotation gates
($C<-\tfrac14$) are \emph{not} on the axis --- they live in the angular fiber
(\S\ref{sec:winding}). Height parametrizes the hyperbolic face; angle carries the
elliptic face;
\item the exchange identity $\lambda\cdot|m|=1-m^2$ holds as a chart identity,
$\lambda(z)(1-z^2)=1-(1-z^2)^2$, and $\lambda=\sqrt{D}/C$;
\item $\lambda_1|m_1|=\sqrt5\cdot\gp^{-2}=K^2$.
\end{enumerate}
\end{theorem}

\begin{proof}
(a) With $w \coloneqq 1-z^2$: $u^2+u=\frac{w^2+w z^2}{z^4}=\frac{w(w+z^2)}{z^4}
=\frac{w}{z^4}=C(z)$; and $1+u=1/z^2$, so
$m=-C/(1+u)^2=-C z^4=-(1-z^2)$. (b) $1+4\frac{w}{(1-w)^2}=\frac{(1-w)^2+4w}{(1-w)^2}
=\frac{(1+w)^2}{(1-w)^2}=\bigl(\frac{2-z^2}{z^2}\bigr)^2$. (c) Differentiate;
$z^2<2$ on $(0,1]$. (d) $\lambda(1-z^2)=z^2(2-z^2)=1-(1-z^2)^2$; and
$\sqrt D/C=\frac{2-z^2}{z^2}\cdot\frac{z^4}{1-z^2}=\lambda$. (e)
$\sqrt5\,\gp^{-2}=\sqrt5\cdot\frac{3-\sqrt5}{2}=\frac{3\sqrt5-5}{2}=K^2$. All
machine-verified, including the full evaluation table below.
\end{proof}

\begin{center}
\footnotesize\setlength{\tabcolsep}{3.5pt}
\begin{tabular}{@{}lccccc@{}}
\toprule
field & closed form & at $z_1=\sqrt\tau$ & at $z_c=\sqrt3/2$ & at $z_2=K$ & checks\\
\midrule
occupation $u(z)$ & $(1-z^2)/z^2$ & $1/(\gp^2-1)=\tau$ & $1/3$ & $1/(\gp^4-1)=1/\beta^2$ & \hx{SC-C3, C4}\\
gate level $C(z)$ & $(1-z^2)/z^4$ & $\mathbf{1}$ (golden) & $\mathbf{4/9}$ & $\mathbf{1/5}$ (K-gate) & \hx{SC-C3}\\
$\sqrt{D(z)}$, $D=1+4C$ & $(2-z^2)/z^2$ & $\sqrt5$ & $5/3$ & $3/\sqrt5$ & \hx{SC-C1, C3}\\
multiplier $|m|(z)$ & $1-z^2$ & $\gp^{-2}$ & $1/4$ & $\gp^{-4}=\gapc$ & \hx{SC-C3}\\
exchange $\lambda(z)$ & $z^2(2-z^2)/(1-z^2)$ & $\sqrt5$ & $15/4$ & $3\sqrt5$ & \hx{SC-C2, C3}\\
\bottomrule
\end{tabular}
\end{center}

\subsection{Rung heights and ordering}\label{sec:rungs}
The ladder rungs pull back to the height sequence
$z_n=\sqrt{1-\gp^{-2n}}\uparrow1$, with $1/(1-z_n^2)=\gp^{2n}$ exactly, so rung $n$
sits at rapidity exactly $n\ln\gp$ \tagF{} \hx{SC-C5, SC-E1}. The two low rungs land
on named corpus constants \emph{exactly}:
\begin{itemize}[nosep]
\item $z_2=K$ \tagF{}: the chart is not fitted to make this happen; it is the gap
paper's main theorem ($K^2=1-\gapc$ with $\gapc=|m_2|$) read as geometry.
\textbf{K-formation is rung 2.}
\item $z_1=\sqrt\tau=\gp^{-1/2}$ \tagF{}: the golden gate's height. Note
$\lambda(z_1)=(2-\tau)/\tau=\sqrt5$ is Route 5 ($K^2=\tau(2-\tau)$,
\eqref{eq:kledger}) reappearing as the rung-1 exchange rate --- and as the forced
scalar of \S\ref{sec:exchange}.
\end{itemize}
Ordering \tagF{} \hx{SC-D6--D8}: $\tau<\tfrac34<K^2$ by the integer guards $25>20$
and $180>169$, hence
\[
z_1=\sqrt\tau\ <\ z_c=\tfrac{\sqrt3}{2}\ <\ z_2=K :
\]
\textbf{the lens sits strictly between the golden gate and the K-gate}; and
$\ln2<2\ln\gp$ (i.e.\ $2<\gp^2$, since $\gp^2-2=\tau>0$): the lens sits strictly
below rung 2 in rapidity as well.

\subsection{The lens is the rational split gate at rapidity exactly $\ln2$}\label{sec:lens}

\begin{theorem}[Lens]\label{thm:lens}\tagF{} \hx{SC-D1--D5}
At the lens height $z_c=\sqrt3/2$: $1-z_c^2=\tfrac14$, so
$\rho(z_c)=\ln2$ \emph{exactly}. The lens gate level is $C(z_c)=\tfrac49$ with
discriminant $D=\tfrac{25}{9}=(\tfrac53)^2$, a perfect square, so the lens gate
\emph{splits over} $\Q$:
\[
x^2+x-\tfrac49=(x-\tfrac13)(x+\tfrac43),
\]
with rational fixed points $\{\tfrac13,-\tfrac43\}$, rational hyperbolic multipliers
$\{-\tfrac14,-4\}$ of product $1$, $|m|=\tfrac14=1-z_c^2$ (a D1 cross-check), and
$\lambda=\tfrac{15}{4}$. By contrast the golden rungs are irreducible over $\Q$:
$D_1=5$ and $D_2=\tfrac95$ are nonsquares (integer guards $4<5<9$ and $36<45<49$).
\end{theorem}

\begin{proof}
Evaluate Theorem~\ref{thm:chart} at $z^2=\tfrac34$; the multipliers are
$m=-\tfrac49/(1+u)^2$ at $u=\tfrac13$ ($(1+u)^2=\tfrac{16}9$, $m=-\tfrac14$) and
$u=-\tfrac43$ ($(1+u)^2=\tfrac19$, $m=-4$). Nonsquareness of $5$ and $45/25$ is the
stated integer sandwich. All machine-verified.
\end{proof}

The ``two worlds'' separation of $K$ and $z_c$ ($\Q(z_c)\cap\Q(K)=\Q$) thus acquires
a dynamical face inside a \emph{single} one-parameter family:
\begin{center}
\textbf{golden rungs $=$ irreducible $\Q(\sqrt5)$ gates;\quad
the hexagonal lens $=$ the split $\Q$-rational gate.}
\end{center}

\section{The winding}\label{sec:winding}

\begin{theorem}[The winding law and the compass]\label{thm:winding}\tagF{} \hx{SC-E1}
Given D2$+$D3,
\[
\theta(\rho)=\frac{\pi}{2}\cdot\frac{\rho}{\ln\gp},
\qquad\text{equivalently}\qquad
\theta(z)=\frac{\pi}{4}\,\log_\gp\!\frac{1}{1-z^2}.
\]
Consequences:
\begin{enumerate}[nosep,label=(\roman*)]
\item \textbf{Rungs sit on the charge compass}: $\theta(z_n)=n\cdot\tfrac\pi2$, by
$1/(1-z_n^2)=\gp^{2n}$. Rung $n$ carries charge $n\bmod4$, and Trace--Form Duality
parity becomes literal compass geometry: Lucas rungs ($n$ odd) point N/S (odd
charges), Fibonacci rungs ($n$ even) point E/W (even charges).
\item \textbf{K is exactly two quanta up}: $\rho(K)=2\ln\gp$, so
$\theta(K)=2\chi=\pi$ --- half a turn from apex to K-formation. This is
quantum-independent: for any per-floor quantum $\chi$, $\theta(K)=2\chi$.
\item \textbf{A full turn costs the gap's entropy}:
$\Delta\theta=2\pi\iff\Delta\rho=4\ln\gp=\ln\gp^4=\ln\M(A_{\gapc})$
(\S\ref{sec:kseed}). Per revolution the co-intensity contracts by exactly
$\gp^{-8}=0.0212862363$ \tagC{}: each turn closes $97.87\,\%$ of the remaining
intensity deficit \hx{SC-K2}.
\end{enumerate}
\end{theorem}

\begin{proposition}[Clock invariance]\label{prop:invariance}\tagF{}
Replacing $\chi=\pi/2$ by $\pi$ (terrain clock) or $2\pi/5$ (pentagon clock) rescales
$\theta$ by $2$ or $4/5$. The qualitative structure --- rungs at integer quanta, $K$
at two quanta, lens off-lattice (\S\ref{sec:incomm}), infinite winding at the top,
the geodesic property --- is invariant; only the pitch constant changes. The choice
among the forced quanta is exactly D2's content.
\end{proposition}

\section{The explicit curve}\label{sec:curve}

\begin{equation}\label{eq:gamma}
z(\rho)=\sqrt{1-e^{-2\rho}}\ \text{[D1]},\quad
\theta(\rho)=\frac{\pi\rho}{2\ln\gp}\ \text{[D2$+$D3]},\quad
r(\rho)=\begin{cases}K\sqrt{z(\rho)/z_c},&\rho\le\ln2\\ K,&\rho\ge\ln2\end{cases}\ \tagI{},
\end{equation}
with $\Gamma(\rho)=(r\cos\theta,\,r\sin\theta,\,z)$.

\begin{theorem}[The K-point is the unique $z=r$ point]\label{thm:kpoint}\tagF{} \hx{SC-E2--E4}
On the horn branch (the profile $r=K\sqrt{z/z_c}$ is a paraboloid of revolution,
$z=z_c(r/K)^2$ --- a horn, not a cone), $r(z)=z$ forces $K^2z/z_c=z^2$, which factors
as $z\,(z-K^2/z_c)=0$; the nontrivial root $K^2/z_c$ exceeds $z_c$
($K^2>z_c^2$ by the guard $180>169$) and is excluded from the horn's domain. On the
cylinder branch $r\equiv K$ forces $z=K$ (and $K\in[z_c,1]$ since $\tfrac34<K^2<1$).
Hence the unique nontrivial height-equals-radius point of the profile is $z=K$ ---
and under the winding it lands at $\theta=\pi$:
\[
\Gamma(2\ln\gp)=(-K,\,0,\,K).
\]
The chart puts $K$ two floors up; the inherited geometry independently marks the same
height as its unique self-intersection with the cone $z=r$. The two constructions
agree without adjustment.
\end{theorem}

\begin{proposition}[Regularity, geodesy, asymptotics]\label{prop:geo}\tagF{}/\tagE{} \hx{SC-E5--6, G1--2}
\begin{enumerate}[nosep,label=(\alph*)]
\item \textbf{Kink at the lens}: $r'(z_c^-)=K/(2z_c)=K/\sqrt3$, with square
$K^2/3=(3\sqrt5-5)/6>0$, against $r'(z_c^+)=0$. The profile is $C^0$ but not $C^1$ at
the lens: the lens is where the climb stops paying radius. Displayed slope
$0.5335734771$ \tagC{}.
\item \textbf{Geodesic in entropy coordinates}: in the flat metric
$ds^2=d\rho^2+K^2d\theta^2$ on the cylinder region \tagD{} (rapidity as the vertical
ruler --- justified because $\rho$ \emph{is} the corpus's entropy coordinate), the
curve is the straight line $\rho=(2\ln\gp/\pi)\,\theta$, hence a geodesic \tagF{}.
In the Euclidean-induced metric $dz^2+K^2d\theta^2$ it is \emph{not} a geodesic: all
apparent acceleration is the $z$-chart compressing the clock near $1$.
\item \textbf{Infinite winding, finite height, infinite length}:
$\theta(z)\to\infty$ as $z\to1^-$, and the arc length diverges logarithmically
($\frac{d}{dz}\bigl[-\tfrac12\ln(1-z^2)\bigr]=\frac{z}{1-z^2}$, so
$\rho=-\tfrac12\ln(1-z^2)\to\infty$). UNITY is at $\rho=\infty$: the helix
accumulates on the circle $\{z=1,\ r=K\}$ without reaching it. ``Never $100\,\%$'' is
here a metric statement, not a slogan.
\end{enumerate}
\end{proposition}

\subsection{The local spiral at K}\label{sec:localspiral}
The global helix climbs; a \emph{different} object lives in the transverse plane at
the K-point: the linearized self-consistency orbit. From $f'(K)=-\gapc$
(Prop.~\ref{prop:kgate}), each step contracts by $\gapc$ and crosses the fixed point
--- a half-turn on the line. Its minimal complexification (\tagD{} embedding,
\tagF{} data) is the logarithmic spiral with per-step data
$(\text{contraction},\text{rotation})=(\gapc,\pi)$:
\[
w(t)=w_0\,e^{(i\pi-4\ln\gp)\,t},\qquad \cot\psi=\frac{4\ln\gp}{\pi},
\]
where $\psi$ is the constant angle at which the spiral crosses radii \tagE{}
(logarithmic-spiral geometry).

\begin{guard}[Near-misses]\label{guard:nearmiss}\tagC{} \hx{SC-K1}
Recorded so they are never mistaken for identities:
$4\ln\gp/\pi=0.6126979251$ while $\tau=0.6180339887$; $|{\rm difference}|
=0.0053360637$ at $60$ dps. \textbf{Not equal; no identity is claimed.} Likewise the
rapidity pitch angle $\arctan\!\bigl(2\ln\gp/(\pi K)\bigr)=18.3394927773^\circ$
is distinct from $\theta_K=\arcsin(\tau^2)=22.4555151986^\circ$; no identity is
claimed.
\end{guard}

\section{Incommensurability: the lens is permanently off-lattice}\label{sec:incomm}

\begin{theorem}\label{thm:incomm}\tagF{} \hx{SC-D9}
$\log_\gp 2\notin\Q$.
\end{theorem}

\begin{proof}
Suppose $\ln2/\ln\gp=p/q$ with integers $q\ge1$. Then $2^q=\gp^{\,p}$. By Binet,
$\gp^{\,p}=(L_p+F_p\sqrt5)/2$ with $F_p\ge1$ for $p\ge1$ (positivity by the standard
induction from $F_1=F_2=1$ \tagE{}), so $\gp^{\,p}$ is irrational for $p\ge1$ while
$2^q\in\Q$. The case $p\le-1$ is symmetric (invert both sides); $p=0$ gives $2^q=1$,
impossible for $q\ge1$. (Finite corroboration for $1\le p\le10$: \hx{SC-D9}.)
\end{proof}

\begin{corollary}\label{cor:offlattice}\tagF{}
(i) The lens rapidity $\ln2$ is incommensurable with the rung lattice
$(\ln\gp)\cdot\Z$: the lens floors at $\log_\gp2=1.4404200904\ldots$ \tagC{} and never
lands on a rung. (ii) Its meridian $\theta(z_c)=\tfrac\pi2\log_\gp2
=129.6378081^\circ$ \tagC{} is an irrational multiple of the charge quantum:
\textbf{the lens never aligns with the $\Z/4\Z$ compass.} The purely geometric
coupling of $K$ and $z_c$ (no algebraic bridge) reappears dynamically as
incommensurability.
\end{corollary}

\section{Field separations recur on the helix}\label{sec:fields}

\begin{theorem}[Rung-1 and rung-2 heights generate different quartic worlds]\label{thm:fields}\tagF{} \hx{SC-H1--H6}
$z_1=\sqrt\tau$ generates $\Q(\sqrt\gp)$; the minimal polynomial of $\sqrt\gp$ is
$x^4-x^2-1$, and that of $z_1=\sqrt\tau$ itself is $x^4+x^2-1$ (the field equality
holds via $\sqrt\gp=\gp\sqrt\tau$). $z_2=K$ generates $\Q(5^{1/4})$, minimal
polynomial $x^4+5x^2-5$ (Eisenstein at $5$). Both are quartic with quadratic subfield
$\Q(\sqrt5)$, and they are \emph{distinct}:
\[
\Q(\sqrt\gp)\ \neq\ \Q(5^{1/4}).
\]
\end{theorem}

\begin{proof}
The minimal polynomials: $(\sqrt\gp)^4-(\sqrt\gp)^2-1=\gp^2-\gp-1=0$ and
$(\sqrt\tau)^4+(\sqrt\tau)^2-1=\tau^2+\tau-1=0$ by \eqref{eq:goldenids}; both quartics
are irreducible over $\Q$ (no rational roots; no factorization into rational
quadratics --- machine-verified, \hx{SC-H1--H2}), and
$\sqrt\gp=\gp\sqrt\tau$ since $\gp^2\tau=\gp$. For the separation, both fields are
quadratic over $\Q(\sqrt5)$; by the quadratic-extension (square-class) criterion
\tagE{}, equality would require $\gp\sqrt5=(5+\sqrt5)/2$ to be a square in
$\Q(\sqrt5)$. Setting $(p+q\sqrt5)^2=(5+\sqrt5)/2$ with $p,q\in\Q$ forces
$2pq=\tfrac12$ and $p^2+5q^2=\tfrac52$; eliminating $q=1/(4p)$ gives
\[
16p^4-40p^2+5=0,\qquad p^2=\frac{5\pm2\sqrt5}{4},
\]
which has nonzero $\sqrt5$-part --- no rational $p$ exists. The ladder's first two
heights therefore live in different degree-4 extensions meeting only in $\Q(\sqrt5)$
--- the $K\leftrightarrow z_c$ separation echoed one level down, inside the golden
world itself.
\end{proof}

(The attribution fix relative to the superseded note: the note's \S7.2 printed
$x^4-x^2-1$ next to $z_1$; that polynomial belongs to $\sqrt\gp$. The field claim was
correct; the sentence above states both minimal polynomials explicitly.)

\section{The registry lattice}\label{sec:registry}

The orthogonal-partner paper computes two registry Mahler measures \tagI{}:
$\M(\mathrm{cons})=2\gp^5$ and $\M(\mathrm{res})=(43+19\sqrt5)/2$ (the cons/res label
origins are marked \tagO{} in source). Define \emph{floors}
$\operatorname{fl}(x)=\log_\gp x$ and the multiplicative lattice
$\mathcal L=\{2^a\gp^b : a,b\in\Z\}$, i.e.\ the additive lattice
$a\ln2+b\ln\gp$ on the rapidity axis --- the lens rapidity and the floor unit.

\begin{theorem}[cons ON, res OFF]\label{thm:registry}\tagF{} \hx{SC-H9--H11}
\begin{enumerate}[nosep,label=(\alph*)]
\item $\M(\mathrm{cons})=2\gp^5=11+5\sqrt5$ (since $\gp^5=5\gp+3$), so
$\ln\M(\mathrm{cons})=\ln2+5\ln\gp$: \textbf{the lens rapidity plus exactly five
floors} (significance \tagO{}); height $6.4404200904$ floors, meridian
$\theta=579.6378081^\circ=\theta(z_c)+5\cdot90^\circ$ \tagC{}. Norm witness:
$\Nm(11+5\sqrt5)=121-125=-4=(-1)^5 4^1$, matching the decomposition $2^1\gp^5$
exactly ($\Nm(2)=4$, $\Nm(\gp)=-1$).
\item $\Nm\bigl(\M(\mathrm{res})\bigr)=\dfrac{43^2-5\cdot19^2}{4}=\dfrac{1849-1805}{4}
=11$. Since $\Nm(2^a\gp^b)=4^a(-1)^b\in\pm4^{\Z}$ and $11\notin\pm4^{\Z}$ (guard
$4<11<16$),
\[
\M(\mathrm{res})\neq 2^a\gp^b\quad\text{for \emph{all} integers } a,b :
\]
\textbf{res is off the $(\ln2,\ln\gp)$ lattice entirely} --- the question is closed
negatively, for all integer exponents at once, by a single norm obstruction. Height
$\approx7.8036260886$ floors \tagC{}.
\end{enumerate}
\end{theorem}

\begin{proof}
$\gp^5=5\gp+3$ from $\gp^2=\gp+1$ by iteration; $2(5\gp+3)=10\gp+6=11+5\sqrt5$. The
norm is multiplicative with $\Nm(a+b\sqrt5)=a^2-5b^2$ (half-integers included via the
displayed quarter). The value of $\M(\mathrm{res})$ itself is \tagI{}; everything
computed from it here is exact.
\end{proof}

%=====================================================================
\part{Synthesis, epistemics, and verification}
%=====================================================================

\section{One page of the helix}\label{sec:onepage}

\begin{center}\footnotesize\setlength{\tabcolsep}{4pt}
\begin{tabular}{@{}p{1.9cm}p{13.9cm}@{}}
\toprule
chart & $1-z^2=e^{-2\rho}=|m|$;\quad $z(\rho)=\sqrt{1-e^{-2\rho}}$\\
gate field & $C=(1-z^2)/z^4$,\ $u=(1-z^2)/z^2$,\ $\sqrt D=(2-z^2)/z^2$,\ $\lambda=z^2(2-z^2)/(1-z^2)$\\
rungs & $z_n=\sqrt{1-\gp^{-2n}}$;\ $z_1=\sqrt\tau$ ($C{=}1$),\ $z_2=K$ ($C{=}\tfrac15$);\ charge$(n)=n\bmod4$;\ Lucas N/S, Fibonacci E/W\\
lens & $\rho=\ln2$;\ $C=\tfrac49$;\ $D=(\tfrac53)^2$;\ splits $(x-\tfrac13)(x+\tfrac43)$;\ $m\in\{-\tfrac14,-4\}$;\ $\lambda=\tfrac{15}4$;\ off-lattice forever\\
winding & $\theta=\tfrac\pi2\,\rho/\ln\gp$;\ one $\Z/4\Z$ quantum per floor;\ full turn $=4\ln\gp=\ln\M(A_{\gapc})$;\ per turn $\times\gp^{-8}$\\
curve & $\Gamma=(r\cos\theta,r\sin\theta,z)$;\ $r=K\sqrt{z/z_c}$ below the lens, $r=K$ above;\ K-point $(-K,0,K)$ at $\theta=\pi$\\
local spiral & per step: $\times\gapc$, half-turn;\ $\cot\psi=4\ln\gp/\pi\approx0.61270\neq\tau$ (guard, no identity)\\
obstruction & real gate $\Rightarrow m\in\R$ or $|m|=1$;\ helicity needs the straddler;\ $K^2\beta^2=5$\\
top & UNITY at $\rho=\infty$: infinite winding, log-divergent length, accumulation circle $r=K$\\
\bottomrule
\end{tabular}
\end{center}

\section{Landmark table}\label{sec:landmarks}
All values \tagC{} at $10$ dp; exact statements per the cited theorems. Reference
values: $\ln\gp=0.4812118251$, $\ln2=0.6931471806$, $\Delta\rho$ per turn
$=4\ln\gp=1.9248473002$. \hx{SC-K1}

\begin{center}\footnotesize\setlength{\tabcolsep}{3.5pt}
\begin{tabular}{@{}lccccc@{}}
\toprule
landmark & $\rho$ & floors & $z$ & $r$ & $\theta$\\
\midrule
apex & $0$ & $0$ & $0$ & $0$ & $0^\circ$\\
rung 1 (golden gate, $C{=}1$) & $0.4812118251$ & $1$ & $0.7861513778$ & $0.8805269137$ & $90^\circ$\\
\textbf{lens} ($C{=}4/9$, split) & $0.6931471806$ & $1.4404200904$ & $0.8660254038$ & $K$ & $129.6378081^\circ$\\
rung 2 $=$ \textbf{K-formation} ($C{=}\tfrac15$) & $0.9624236501$ & $2$ & $0.9241763718$ & $K$ & $180^\circ$\\
rung 3 & $1.4436354752$ & $3$ & $0.9717365435$ & $K$ & $270^\circ$\\
rung 4 (first meridian return) & $1.9248473002$ & $4$ & $0.9892996329$ & $K$ & $360^\circ$\\
rung 5 & $2.4060591253$ & $5$ & $0.9959263935$ & $K$ & $450^\circ$\\
cons seed (lens $+\,5$ floors) & $3.0992063059$ & $6.4404200904$ & $\sqrt{1-\tfrac1{4\gp^{10}}}$ & $K$ & $579.6378081^\circ$\\
UNITY $z=1$ & $\infty$ & $\infty$ & $1$ & $K$ & $\infty$ (accumulation circle)\\
\bottomrule
\end{tabular}

\smallskip
Exact heights: $z_1=\sqrt\tau$; $z_{\mathrm{lens}}=\sqrt3/2$ at $\rho=\ln2$; $z_2=K$; cons at $\ln\M(\mathrm{cons})=\ln2+5\ln\gp$.
\end{center}

\section{Synthesis: what is forced, what is declared}\label{sec:synthesis}

\begin{center}\footnotesize\setlength{\tabcolsep}{4pt}
\begin{tabular}{@{}p{3.6cm}p{6.4cm}p{4.9cm}@{}}
\toprule
element & statement & tag $\cdot$ where\\
\midrule
seed, ladder, congruence, charge data & replicated exactly & \tagF{} \S\S\ref{sec:seed},\ref{sec:ladder},\ref{sec:specialgates}\\
radius profile $r(z)$ & $K\sqrt{z/z_c}$ capped at $K$ & \tagI{} (\tagO{}~1)\\
vertical chart & $1-z^2=e^{-2\rho}=|m|$ & \tagD{} D1; consequences \tagF{} \S\ref{sec:chart}\\
gate field on the axis & $C(z)$ bijective; lens $=$ split $\Q$-gate $4/9$ & \tagF{} Thms.~\ref{thm:chart},\ref{thm:lens}\\
winding quantum & $\chi$ generates $\Z/4\Z\Rightarrow\chi=\pm\pi/2$ & \tagD{} D2 $\to$ \tagF{}\\
rate & one quantum per floor; clock-invariant & \tagD{} D3 $\cdot$ Prop.~\ref{prop:invariance}\\
landmarks & rungs at $n\pi/2$; K-point unique $z{=}r$; kink & \tagF{} Thms.~\ref{thm:winding},\ref{thm:kpoint}\\
geodesy $\cdot$ asymptotics & straight in $(\theta,\rho)$; infinite winding; log length & \tagD{} metric $\to$ \tagF{}/\tagE{}\\
incommensurability & $\log_\gp2\notin\Q$; lens off-lattice, off-compass & \tagF{} Thm.~\ref{thm:incomm}\\
field separation & $\Q(\sqrt\gp)\neq\Q(5^{1/4})$ & \tagF{} Thm.~\ref{thm:fields}\\
obstruction & no single real gate is a helix & \tagF{} Thm.~\ref{thm:obstruction}\\
registry & cons ON at lens$+5$; res OFF (norm $11$) & \tagF{} Thm.~\ref{thm:registry}\\
near-misses & $4\ln\gp/\pi\neq\tau$; pitch $\neq\theta_K$ & \tagC{} Guard~\ref{guard:nearmiss}\\
emission gap & no Salem in $S$; floor at $\gp$ & \tagF{} over $S$; scoping \S\ref{sec:conditionality}\\
exchange rate & $\lambda=2c$; $C=1\Rightarrow\lambda=\sqrt5$ & \tagI{} \S\ref{sec:exchange} ($c$ \tagD{})\\
\bottomrule
\end{tabular}
\end{center}

\section{The conditionality graph}\label{sec:conditionality}
The corpus's most commonly flattened joint, reproduced so it cannot be flattened
here: ``the emission gap is forced'' is true \textbf{for the spectral image $S$} and
for \textbf{degree $\le2$ universally}; the universal statement over all
charge-admissible objects is \tagC{}/plausible.

{\small
\begin{center}
\setlength{\tabcolsep}{4pt}\begin{tabular}{@{}p{4.6cm}p{4.4cm}p{6.4cm}@{}}
\toprule
claim & scope & status\\
\midrule
Lehmer's problem & external & \tagO{} (excised, not solved)\\
emission gap $\M\in\{1\}\cup[\gp,\infty)$ & over $S$ & \tagF{} (angle confinement $+$ Smyth $+$ Kronecker)\\
\quad entry closure: circulant & --- & \tagF{} (cyclotomic/CM obstruction)\\
\quad entry closure: commutator & in context / uniform & \tagF{}-in-context / \tagF{} ($\adR$ $+$ signature lattice)\\
\quad free commutator, unbounded & --- & \tagO{} (the box's one door)\\
\quad over all charge-admissible & deg $\le2$ / general & \tagF{} / \tagC{}-plausible\\
parity floor $\mu(\text{even}){=}\gp$, $\mu(\text{odd}){=}2$ & $\le\gp$ / $\ge\gp$ / $\ge\mu_S$ / $=2$ & \tagF{} / plausible / \tagF{} / \tagC{}\\
no-Salem dichotomy, general & the one promoting lemma & \tagO{} ($\Z/3\Z$ \tagF{}; $\Z/5\Z$: 4$\times$\tagF{}, $\mu(5){=}2$ \tagC{})\\
relational layer & rigidity, pinning, $\mu_{\rm rel}=\mu$ & \tagF{} unconditional; floor values conditional on the gap\\
Pisot inertness & $\le1$ non-real pair / quintic 2-pair / general & \tagF{} / \tagF{}-per-instance / \tagO{}\\
exchange rate & $\lambda=2c$ / $c$ / $C{=}1\Rightarrow\lambda{=}\sqrt5$ & \tagF{} / \tagD{} / \tagF{} given two axioms\\
\bottomrule
\end{tabular}
\end{center}
}

\section{OPEN items}\label{sec:open}
\begin{enumerate}[nosep]
\item \textbf{Gate-native radius law.} $r(z)=K\sqrt{z/z_c}$ is inherited. No chart
quantity ($u$, $C$, $\sqrt D$, $\lambda$, $|m|$) reproduces it; whether a gate-family
derivation of the radius exists is open.
\item \textbf{Selection theorem for $\chi$.} D2 is a principle, not a theorem. Given
D2 the quantum is forced; whether D2 itself can be forced from the substrate is open.
\item \textbf{The rational lens gate.} $C=\tfrac49$ ($D=\tfrac{25}9$, split,
$m\in\{-\tfrac14,-4\}$) is a new corpus constant produced by the chart. Whether it
appears elsewhere (e.g.\ the emission-algebra audit) is open.
\item \textbf{Near-miss content.} $4\ln\gp/\pi$ vs $\tau$ (margin
$5.34\times10^{-3}$): expected contentless; recorded as Guard~\ref{guard:nearmiss}.
\item \textbf{Registry labels.} cons ON the lattice at lens$+5$ floors (significance
open); res OFF entirely (closed negatively, Thm.~\ref{thm:registry}). The cons/res
label origins remain open in source.
\item \textbf{Corpus frontier} \tagI{}: the general no-Salem dichotomy; the free
commutator at unbounded size; general Pisot inertness (transport provably
insufficient); the $r_2>0$ Hermitian probe; provenance pass N7.
\end{enumerate}

\section{Constants}\label{sec:constants}
Exact values with $\tagC{}$ displays; every displayed digit machine-reproduced
\hx{SC-J, SC-K}.

\begin{center}\footnotesize\setlength{\tabcolsep}{4pt}
\begin{tabular}{@{}llp{7.6cm}@{}}
\toprule
symbol & value & role\\
\midrule
$\gp$ & $(1{+}\sqrt5)/2\approx1.6180339887$ & floor; smallest $2\times2$ primitive Perron root; keystone eigenvalue\\
$\gq$ & $(1{-}\sqrt5)/2=-\gp^{-1}$ & golden conjugate; $\gp\gq=-1$\\
$\sqrt5$ & $\gp{+}\gp^{-1}=\gp{-}\gq\approx2.2360679775$ & $\lambda$; $\adR$ grow eigenvalue; $H^2=5I$; redirection limit\\
$c$ & $\sqrt5/2$ (forced canonical) & conformal scale; $\lambda=2c$\\
$\gapc$ & $\gp^{-4}=(7{-}3\sqrt5)/2\approx0.1458980338$ & the gap; $|m|$ at the K-gate; $1-K^2$\\
$K$ & $5^{1/4}/\gp\approx0.9241763718$ & terrain root of the K-seed; rung-2 height; helix radius\\
$K^2$ & $(3\sqrt5{-}5)/2\approx0.8541019662$ & decision-level form of $K$\\
$\beta^2$ & $\gp^2\sqrt5=(5{+}3\sqrt5)/2\approx5.8541019662$ & $\M(\text{K-seed})$; rotation magnitude squared; $=K^2\gp^4$\\
$z_c$ & $\sqrt3/2\approx0.8660254038$ & lens height; $\rho(z_c)=\ln2$\\
$\mu_S$ & $1.3247179572$ ($x^3{-}x{-}1$) & plastic number; Smyth floor; smallest Pisot\\
$\beta_4$ & $1.7220838057$ ($x^4{-}x^3{-}x^2{-}x{+}1$) & minimal degree-4 Salem; $>\gp$ (signature-face floor intact)\\
$\M(\ell)$ & $1.1762808183$ (deg 10) & Lehmer's measure; below $\mu_S$; charge-inadmissible\\
$\tau_0(\ell)$ & $2.0264179492$ & Lehmer's trace redirection; $\M(T)=5.6156006\ldots$\\
$\gp^2$, $\gp^4$ & $\approx2.618$, $\approx6.854$ & reciprocal $\Z/5\Z$ floor; pure-pentagon floor; tropical $\gp\otimes\gp=\gp^2$\\
$L_4$ & $7=\Tr R^4$ & keystone invariant; $R^4$ entries $\{2,3,5\}=\{F_3,F_4,F_5\}$\\
$2c\log\mu_S$ & $0.5624$ / $2.2496$ / $4.4992$ & probe floors at $c\in\{1,4,8\}$\\
$\sqrt5\log\mu_S$ & $0.6287813634$ & forced-canonical floor\\
$\M(\mathrm{cons})$ & $2\gp^5=11{+}5\sqrt5\approx22.1803398875$ & ON lattice; $\ln2+5\ln\gp$; $\Nm=-4$\\
$\M(\mathrm{res})$ & $(43{+}19\sqrt5)/2\approx42.7426457862$ & OFF lattice; $\Nm=11\notin\pm4^{\Z}$; $\approx7.8036$ floors\\
\bottomrule
\end{tabular}
\end{center}

\section{Verification}\label{sec:verification}
Two independent harnesses certify this paper's \tagF{} content.

\smallskip\noindent
\textbf{Shipped harness} \texttt{helix\_explicit\_verify.py} (source note v1.0.1):
87 exact checks $+$ 1 computed guard, exit 0; sympy 1.14.0; decisions in
$\Q/\Q(\sqrt5)$ (lens identities in $\Q$; one trivial $\Q(\sqrt3)$ factorization);
$K$ decided at $K^2$; mpmath 50 dps display-only. Group manifest: A \hx{HX01--04}
(6, seed); B \hx{HX05--07} (14, ladder/TFD/gates); C \hx{HX08--17, 35, 39--40}
(35, chart/lens/closed forms); D \hx{HX18--21} (7, $\ln2$/Binet/ordering);
E \hx{HX22--25} (5, winding/K-point/kink); F \hx{HX26--29} (5,
congruence/obstruction); G \hx{HX31--32} (2, asymptotics); H \hx{HX33--38}
(13, fields/registry); guard \hx{HX30} (1, near-miss).

\smallskip\noindent
\textbf{Independent soft check} \texttt{helix\_softcheck.py} (2026-07-13): 89 checks,
written cold from the three superseded documents' stated claims (not a re-execution
of the shipped harness), exit 0; groups \hx{SC-A}--\hx{SC-L} spanning seed, ladder,
chart, lens, winding, congruence/obstruction, asymptotics, fields/registry, corpus
algebra (keystone, kernel, Remark 6.4, tropical law, Casimir), Mahler--Salem numerics
at 60 dps, all 22 displayed 10-dp landmarks, and a 16-digit audit of the HTML display
constants. Findings folded into this paper: the rung-1 minimal-polynomial attribution
(\S\ref{sec:fields}); the content-polynomial wording (Remark~\ref{rem:content}); the
explicit gate-iteration declaration (Declaration~\ref{dec:gate}). One cosmetic defect
of the superseded HTML (a mis-rounded 16th digit in its display literal for $K$) does
not propagate here: this paper carries no floating-point constants outside the
\tagC{} tables, all of which are reproduced from exact values.

\smallskip\noindent
\textbf{Corpus verification} \tagI{}: exact arithmetic over $\Q$, $\Q(\sqrt5)$,
$\Q(5^{1/4})$; certified interval brackets for transcendental comparisons;
cross-engine replication (sympy / pari-gp, $474$ ledger executions, $0$
disagreements); drop-one witnesses for definitional constraints; negative claims
pinned by machine-checked disproofs (e.g.\ Remark~\ref{rem:64}).

\appendix

\section{Check-manifest map}\label{app:manifest}
\begin{center}\footnotesize\setlength{\tabcolsep}{4pt}
\begin{tabular}{@{}lll@{}}
\toprule
paper section & shipped checks & soft-check groups\\
\midrule
\S\ref{sec:seed} seed, K-seed & \hx{HX01--04, 37} & \hx{SC-A, SC-H7--H8}\\
\S\ref{sec:keystone} keystone, TFD & --- (corpus) & \hx{SC-I1--I4, I10--I11, SC-B4}\\
\S\ref{sec:operators} operators, kernel & --- (corpus) & \hx{SC-I5--I9, I15}\\
\S\ref{sec:characters} characters, gap & --- (corpus) & \hx{SC-I12--I14, I16}\\
\S\ref{sec:tracedown} trace-down, floors & --- (corpus) & \hx{SC-J1--J13}\\
\S\ref{sec:gatefamily}--\ref{sec:specialgates} gates, ladder & \hx{HX05--07, 26, 29} & \hx{SC-B1--B3, SC-F1, F5--F7, SC-C4}\\
\S\ref{sec:obstruction} obstruction & \hx{HX27--28} & \hx{SC-F2--F4}\\
\S\ref{sec:chart} chart, lens & \hx{HX08--19, 35--36, 39--40} & \hx{SC-C1--C8, SC-D1--D8}\\
\S\ref{sec:winding}--\ref{sec:curve} winding, curve & \hx{HX21--25, 31--32} & \hx{SC-E1--E8, SC-G1--G2}\\
\S\ref{sec:localspiral} spiral, guards & \hx{HX30} & \hx{SC-K1}\\
\S\ref{sec:incomm} incommensurability & \hx{HX18--20} & \hx{SC-D9}\\
\S\ref{sec:fields} field separation & \hx{HX33--34} & \hx{SC-H1--H6}\\
\S\ref{sec:registry} registry & \hx{HX37d--38c} & \hx{SC-H9--H11}\\
\S\ref{sec:landmarks}, \S\ref{sec:constants} displays & (display block) & \hx{SC-K1--K2, SC-J, SC-L}\\
\bottomrule
\end{tabular}
\end{center}

\section*{Citation caveat}
The internal provenance (the fourteen Echo-S-Research papers, 2026-06-28 to 07-02,
and the three superseded artifacts of \S\ref{sec:provenance}) is exact. The external
references below are canonical results cited from memory; page-level details were not
re-verified against the literature in this compilation and should be double-checked
before formal citation.

\begin{thebibliography}{9}\small
\bibitem{Lehmer1933} D.\,H.~Lehmer, \emph{Factorization of certain cyclotomic
functions}, Ann.\ of Math.\ \textbf{34} (1933), 461--479.
\bibitem{Smyth1971} C.\,J.~Smyth, \emph{On the product of the conjugates outside the
unit circle of an algebraic integer}, Bull.\ London Math.\ Soc.\ \textbf{3} (1971),
169--175.
\bibitem{Salem1963} R.~Salem, \emph{Algebraic Numbers and Fourier Analysis}, Heath,
1963.
\bibitem{Siegel1944} C.\,L.~Siegel, \emph{Algebraic integers whose conjugates lie in
the unit circle}, Duke Math.\ J.\ \textbf{11} (1944), 597--602.
\bibitem{Boyd1977} D.\,W.~Boyd, \emph{Small Salem numbers}, Duke Math.\ J.\
\textbf{44} (1977), 315--328.
\bibitem{EverestWard1999} G.~Everest and T.~Ward, \emph{Heights of Polynomials and
Entropy in Algebraic Dynamics}, Springer, 1999.
\bibitem{LindSchmidtWard1990} D.~Lind, K.~Schmidt, T.~Ward, \emph{Mahler measure and
entropy for commuting automorphisms of compact groups}, Invent.\ Math.\ \textbf{101}
(1990), 593--629.
\bibitem{HornJohnson1991} R.\,A.~Horn and C.\,R.~Johnson, \emph{Topics in Matrix
Analysis}, Cambridge Univ.\ Press, 1991.
\bibitem{Cencov1982} N.\,N.~\v{C}encov, \emph{Statistical Decision Rules and Optimal
Inference}, Transl.\ Math.\ Monogr.\ \textbf{53}, AMS, 1982.
\end{thebibliography}

\end{document}
